\section{Две скорости и мост $\speedC=\caKappaGeom \caSignalSpeed$}
\label{sec:two-c}

На гексагональном срезе $(2{+}1)$ коэффициент
\[
\caKappaGeom_{\mathrm{hex}}=\sqrt{3}/2
\]
(\cref{sec:hex-slice}); на \latGCK
\[
\caKappaGeom_{\mathrm{FCC}}=1/\sqrt{2}
\]
(\cref{sec:fcc-kappa}; индекс $\mathrm{FCC}$ --- тип Браве из \cref{def:packing-types}).
Величина $\caKappaGeom$ --- отношение радиусов вписанной и описанной сфер однотактового тела.
Тождество $\speedC=\caKappaGeom \caSignalSpeed$ служит проверкой согласованности макроописания
после фиксации $\htick=\caKappaGeom t_P$ (\cref{fig:carrier-two-speeds-micro,fig:carrier-two-speeds-macro}).

\begin{definition}[Две скорости и continuum-упаковка $t_P$]
\label{def:two-speeds-formal}
\begin{itemize}
  \item \textbf{Скорость звена на $\Mlayer$:}
  $\caSignalSpeed:=\caStepSpace/\htick$ (\cref{def:minkowski-carrier-split}).
  \item \textbf{Макроскопическая скорость на $\Tlayer$:}
  $\speedC:=\caKappaGeom\,\caSignalSpeed$ (с якорем CODATA после моста к СИ, \cref{def:si-bridge}).
  \item \textbf{Continuum-упаковка:}
  определение $t_P:=\ell_P/\speedC$ в рамке одной скорости $\speedC$ на всех масштабах
  и подстановка $t_P$ вместо $\htick$ при чтении $\caSignalSpeed=\ell_P/t_P$.
\end{itemize}
\end{definition}

\begin{proposition}[Каузальный предел одного такта на $\Mlayer$]
\label{prop:m-one-tick-cone}
За один такт $t\to t+1$ закон~$g$ не переносит информацию за пределы каузальной звезды $N(x)$
(аксиома~$A_1$, \cref{ax:A1}): для дискретного фронта
$\|\Delta x\|_\infty\le \caStepSpace$.
Следовательно, максимальная скорость сигнала на $\Mlayer$ за такт равна $\caSignalSpeed$;
превышение за один $\htick$ без нарушения $A_1$ невозможно.
\end{proposition}

\begin{proposition}[Отношение $\caSignalSpeed/\speedC$ при FCC]
\label{prop:c0-over-c}
При $\caKappaGeom_{\mathrm{FCC}}=1/\sqrt{2}$ и \cref{def:si-bridge}:
\[
  \htick=\caKappaGeom\, t_P,\qquad
  \speedC=\caKappaGeom\,\caSignalSpeed,\qquad
  \frac{\caSignalSpeed}{\speedC}=\frac{1}{\caKappaGeom}=\sqrt{2}.
\]
Величина $\caKappaGeom$ задаётся геометрией однотактового тела (\cref{sec:fcc-kappa}),
не отношением $\speedC/\caSignalSpeed$ как свободным параметром.
\end{proposition}

\begin{proposition}[Предел скорости наблюдаемого фронта на $\Tlayer$]
\label{prop:t-speed-limit}
Скорость распространения EM-вакуумного фронта, извлекаемая прибором на $\Tlayer$
после readout $\mathcal{B}$ на масштабах $\gg\caStepSpace$,
согласована с $\speedC$ (\cref{sec:c-light-cone,ch:macro}; сверка --- CE-M-01).
Утверждение $\caSignalSpeed>\speedC$ в SI \emph{не} определяет на $\Tlayer$
состояние с групповой скоростью $v>\speedC$:
$\caSignalSpeed$ относится к паре $(\caStepSpace,\htick)$ на $\Mlayer$,
$\speedC$ --- к изотропному макро-якорю $\caKappaGeom\,\caSignalSpeed$.
\end{proposition}

\begin{corollary}[Нет кинематического доступа к $\Mlayer$]
\label{cor:no-m-kinematic-access}
По \cref{cor:no-continuum,sec:no-continuum-carrier} шаг $\caStepSpace=\ell_P$ не является
пределом $h\to 0$ для масштаба тела или прибора на $\Tlayer$.
Отображение «уменьшение линейного размера $\Rightarrow$ замена $\speedC$ на $\caSignalSpeed$''
не входит в класс допустимых морфизмов модели: наблюдатель остаётся на $\Tlayer$
с конусом $\speedC$; кинематика $\caSignalSpeed$ --- свойство обновления на $\carrierLattice\times\mathbb{Z}$.
\end{corollary}

\begin{theorem}[Операторный закон квантовой запутанности через отношение скоростей $\caSignalSpeed/\speedC$ и геометрическую лестницу Белла--Цирельсона]
\label{thm:entanglement-c0-tsirelson}
Для любого $d\in\mathbb{N}$ операторный закон Каталана $\mathrm{CatalanLaw}(P,d)$ (\texttt{quantum\_entanglement\_c0\_causal\_holonomy\_and\_tsirelson\_bound} в \texttt{lean/MtCarrier/UnifiedOperators.lean}) эквивалентен системе следующих точных операторных соотношений на оболочке кубооктаэдра $O_2(d)=\mathrm{cubicOrbit3D}(d-1)$:
\begin{enumerate}[label=\textbf{(\roman*)},leftmargin=*]
\item \textbf{Оператор квадрата отношения тактовой скорости звена $\caSignalSpeed$ к скорости света $\speedC$:}
  для каждого вектора смещения $e=(dz,dy,dx)\in O_2(d)$ квадрат евклидовой нормы ребра совпадает с оператором квадрата отношения скоростей $\mathrm{c0OverCSqOp}(d)$ и с отношением числа треугольных граней к числу вершин квадратной грани:
  \begin{equation}
  \label{eq:c0-over-c-sq-operator}
    \forall e\in O_2(d),\qquad
    \|e\|_2^2 \;=\; \mathrm{c0OverCSqOp}(d) \;=\; \caKappaGeom(d)^{-2} \;=\; \left(\frac{\caSignalSpeed}{\speedC}\right)^{\!2} \;=\; \frac{|{\triangle}(d)|}{d_{\mathrm{in}}(\square)}.
  \end{equation}
  Следовательно, фазовая связность вдоль звена кубооктаэдра на микроскопическом $\Mlayer$-слое распространяется со строго конечной тактовой скоростью $\caSignalSpeed = \caKappaGeom(d)^{-1}\speedC$, превышающей макроскопическую скорость электромагнитного фронта $\speedC = \caKappaGeom(d)\caSignalSpeed$ на $\Tlayer$-слое ровно в $\caKappaGeom(d)^{-1}$ раз --- без «мгновенного дальнодействия».
\item \textbf{Операторная геометрическая прогрессия границ Белла--Цирельсона со знаменателем $(\caSignalSpeed/\speedC)^2 = \mathrm{c0OverCSqOp}(d)$:}
  операторы квадратов максимальных корреляторов CHSH для сепарабельной пары $\mathrm{bellSeparableSqOp}(d)$, классического предела локальных скрытых параметров $\mathrm{bellLhvSqOp}(d)$ (ограниченного макроскопическим конусом скорости света $\speedC$ на $\Tlayer$) и квантового предела Цирельсона $\mathrm{tsirelsonQmSqOp}(d)$ (связанного тактовой голономией со скоростью $\caSignalSpeed$ на $\Mlayer$) связаны точными операторными тождествами:
  \begin{equation}
  \label{eq:bell-tsirelson-c0-ladder}
    \mathrm{bellLhvSqOp}(d) \;=\; \mathrm{c0OverCSqOp}(d)\cdot\mathrm{bellSeparableSqOp}(d) \;=\; d_{\mathrm{in}}(\square),
    \qquad
    \mathrm{tsirelsonQmSqOp}(d) \;=\; \mathrm{c0OverCSqOp}(d)\cdot\mathrm{bellLhvSqOp}(d) \;=\; |{\triangle}(d)|.
  \end{equation}
  Тем самым отношение квантового оператора Цирельсона к классическому оператору Белла тождественно равно отношению тактовой скорости звена $\caSignalSpeed$ к скорости света $\speedC$:
  \[
    \frac{\mathrm{tsirelsonQmSqOp}(d)}{\mathrm{bellLhvSqOp}(d)}
    \;=\; \frac{\mathrm{bellLhvSqOp}(d)}{\mathrm{bellSeparableSqOp}(d)}
    \;=\; \frac{|{\triangle}(d)|}{d_{\mathrm{in}}(\square)}
    \;=\; \mathrm{c0OverCSqOp}(d)
    \;=\; \left(\frac{\caSignalSpeed}{\speedC}\right)^{\!2}.
  \]
\item \textbf{Антисимметричная синглетная корреляция Паули при строгом операторном запрете сверхсветовой сигнализации (\emph{no-signaling}):}
  на операторе фазового полупериода Паули $\Phi_{\mathrm{Pauli}}(d)$ чётный оператор голономии антисимметричен относительно вакуумного максимума, а нечётный оператор тока переноса энергии и подпороговый оператор ворот Гейзенберга обращаются в нуль:
  \[
    \mathcal{C}_d\bigl(\Phi_{\mathrm{Pauli}}(d)\bigr) \;=\; -\mathcal{C}_d(0) \;=\; -Q_\square(d),
    \qquad
    \mathcal{S}_d\bigl(\Phi_{\mathrm{Pauli}}(d)\bigr) \;=\; 0,
  \]
  \[
    |\Phi_{\mathrm{raw}}| < \Delta\varphi_{\mathrm{disc}}(d)
    \;\Longrightarrow\;
    \mathrm{heisenbergDeadbandOp}(d,\Phi_{\mathrm{raw}}) = 0
    \;\wedge\;
    \mathrm{energyQuantumOperator}(d,0) = 0.
  \]
  Поэтому по тактовому каналу со скоростью $\caSignalSpeed$ передаётся исключительно безэнергетическая спинорная голономия синглета (защищённая от вакуумной декогеренции мёртвой зоной Гейзенберга), тогда как любой перенос кванта энергии или сигнала ($\mathrm{energyQuantumOperator}(d,\Phi) \ge \mu_{\mathrm{vac}}(d)$) требует надпорогового пробоя $|\Phi|\ge\Delta\varphi_{\mathrm{disc}}(d)$ и ограничен на $\Tlayer$-слое скоростью света~$\speedC$.
\end{enumerate}
\end{theorem}

\begin{remark}[Сверка]
\label{rem:two-c-burden}
\cref{prop:c0-over-c} --- CE-M-00 (\texttt{Kappa\_bottom\_up});
\cref{prop:t-speed-limit} --- CE-M-01 (suite \texttt{t\_macro}).
Ложное тождество «$\speedC$ на шаге $\ell_P$ при $t_P=\ell_P/\speedC$'' есть частный случай
неверной continuum-упаковки из \cref{def:two-speeds-formal}.
\end{remark}

\begin{figure}[htbp]
  \centering
  \input{figures/tikz/carrier-two-speeds-micro.tex}
  \caption{Микро: дискретный шаг по решётке, тактовая скорость $\caSignalSpeed$.}
  \label{fig:carrier-two-speeds-micro}
\end{figure}

\begin{figure}[htbp]
  \centering
  \input{figures/tikz/carrier-two-speeds-macro.tex}
  \caption{Макро: осреднённый фронт, $\speedC=\caKappaGeom \caSignalSpeed$.}
  \label{fig:carrier-two-speeds-macro}
\end{figure}

